\section{Essential automorphisms: basic def\/initions and properties}\label{section2}

\subsection{Background on parabolic geometries and their Weyl structures}\label{section2.1}

Let us begin by recalling the def\/initions of parabolic geometries and their Weyl structures (the latter, introduced by A.~\v{C}ap and J.~Slov\'ak in \cite{CS03}, will be central to our notion of essential parabolic structures). Parabolic geometries are certain types of Cartan geometries, which are very general: given a closed subgroup $P$ of a Lie group $G$, a \textit{Cartan geometry of type $(G,P)$} (or \textit{modelled on the homogeneous space $G/P$}) is given by a principal $P$ bundle $\pi: \Gbdle \rightarrow M$, equipped with a \textit{Cartan connection} $\omega$. That is, $\omega \in \Omega^1(\Gbdle,\g)$ satisf\/ies:
\begin{gather}
R_p^*(\omega) = \mathrm{Ad}\big(p^{-1}\big) \circ \omega, \,\, \mathrm{for} \,\, \mathrm{all} \,\, p \in P; \label{ad-inv}\\
\omega(\tilde{X}) = X, \,\, \mathrm{for} \,\, \mathrm{any} \,\, X \in \p, \, \tilde{X} \,\, \mathrm{its} \,\, \mathrm{fundamental} \,\, \mathrm{vector} \,\, \mathrm{f\/ield} \,\, \mathrm{on} \,\, \Gbdle; \label{fund-vfs}\\
\omega(u) : T_u\Gbdle \rightarrow \g \,\, \mathrm{is} \,\, \mathrm{a} \,\, \mathrm{linear} \,\, \mathrm{isomorphism} \,\, \mathrm{for} \,\, \mathrm{all} \,\,u \in \Gbdle. \label{parallelism}
\end{gather}

A Cartan geometry of type $(G,P)$ is a \emph{parabolic geometry} if $G$ is a real or complex semi-simple Lie group, and $P \subset G$ is a parabolic subgroup as in representation theory~-- at the level of Lie algebras, this means, for $\g$ complex semi-simple, that $\p$ must contain a Borel (maximal solvable) subalgebra of $\g$; for $\g$ real semi-simple, the complexif\/ication $\p(\C)$ must contain a Borel subalgebra of $\g(\C)$. (For a more detailed discussion of the basic properties of parabolic subgroups and parabolic geometries, the reader is referred to \cite{CSbook}. Here we only attempt to cite some of the key facts which are germane to the subsequent text.) In particular, in the parabolic setting the Lie algebra $\g$ of $G$ has an induced $\vert k \vert$-grading for some natural number $k$, so $\g = \g_{-k} \dsum \cdots \dsum \g_k$ with $[\g_i,\g_j] \subseteq \g_{i+j}$ and the subalgebra $\g_- = \g_{-k} \dsum \cdots \dsum \g_{-1}$ is generated by $\g_{-1}$. The Lie algebra of the parabolic subgroup $P$ is the parabolic subalgebra $\p = \g_0 \dsum \cdots \dsum \g_{k}$, which has Levi decomposition $\p = \g_0 \dsum \p_+$ with $\g_0$ reductive and $\p_+ = \g_1 \dsum \cdots \dsum \g_k$ the nilradical of $\p$. At the group level, we have a reductive subgroup $G_0 \subset P$ whose Lie algebra is $\g_0$, and $P \isom G_0 \ltimes P_+$ where $P_+ = \exp (\p_+)$ is a normal, nilpotent subgroup of $P$ globally dif\/feomorphic to $\p_+$ via the exponential map.

The above-stated properties of the parabolic pair $(G,P)$ and their Lie algebras are used to identify the following important geometric structures associated to a parabolic geometry $(\Gbdle \rightarrow M,\omega)$ of type $(G,P)$. The orbit space $\Gbdle_0 := \Gbdle/P_+$ of the $P_+$-action on $\Gbdle$ def\/ines a $G_0$-principal bundle $\pi_0: \Gbdle_0 \rightarrow M$, while by def\/inition we also have a $P_+$-principal bundle $\pi_+: \Gbdle \rightarrow \Gbdle_0$. The f\/iltration of $\g$ by $\mathrm{Ad}(P)$-invariant submodules $\g^i = \g_i \dsum \cdots \dsum \g_k$ descends to a f\/iltration of $\g/\p \isom \g_-$ which is invariant under the quotient representation $\overline{\mathrm{Ad}}: P \rightarrow Gl(\g/\p)$, and thus determines a f\/iltration of the tangent bundle on the base space, $TM = T^{-k}M \supset \cdots \supset T^{-1}M$, via the isomorphism
\[
TM \isom_{\omega} \Gbdle \times_{\overline{\mathrm{Ad}}(P)} \g/\p,
\]
(which holds in general for Cartan geometries) and setting $T^iM \isom_{\omega} \Gbdle \times_{\overline{\mathrm{Ad}}(P)} \g^i/\p$. Furthermore, the Cartan connection $\omega$ descends to $\Gbdle_0$ to identify it as a reduction to $G_0$ of the structure group of the \emph{associated graded tangent bundle} $\mathrm{gr}(TM) = \mathrm{gr}_{-k}(TM) \dsum \cdots \dsum \mathrm{gr}_{-1}(TM)$, where $\mathrm{gr}_i(TM) = T^iM/T^{i+1}M$.

The data $(M,\{T^iM\},\Gbdle_0)$~-- consisting of a smooth manifold $M$, a f\/iltration $\{T^iM\}$ of its tangent bundle which satisf\/ies $\mathrm{rk}(T^iM) = \mathrm{dim}(\g^i/\p)$, and a reduction $\Gbdle_0$ of the structure group of $\mathrm{gr}(TM)$ to $G_0$ --, is called an \emph{infinitesimal flag structure of type} $(\g,P)$. The f\/lag structure is \emph{regular} if the Lie bracket of vector f\/ields on $M$ respects the f\/iltration, i.e.\ $[\Gamma(T^iM),\Gamma(T^jM)] \subset \Gamma(T^{i+j}M)$, and if the alternating bilinear form thus induced on $\mathrm{gr}(TM)$ gives it a point-wise Lie algebra structure isomorphic to $\g_-$. When the f\/lag structure is induced by a parabolic geometry of type $(G,P)$, this regularity assumption can be related to an equivalent regularity condition that the curvature of the Cartan connection $\omega$ have strictly positive homogeneity (cf.~3.1.8 of~\cite{CSbook}). A fundamental theorem of parabolic geometry states that, for any regular inf\/initesimal f\/lag structure of type $(\g,P)$, there exists a regular parabolic geometry of type $(G,P)$ which induces it. This parabolic geometry is uniquely determined up to isomorphism by a normalisation condition on the curvature of the Cartan connection, except for a number of parabolic types $(G,P)$ where the reduction of $\mathrm{gr}(TM)$ to $G_0$ provides no additional information and an extra geometric structure is needed for uniqueness (e.g.\ projective structures, where an additional choice of an equivalence class of connections is needed to f\/ix the structure). In these cases, we will assume that the extra geometric structure is included when we speak of the ``regular inf\/initesimal f\/lag structure''. We thus identify the geometric structure of an inf\/initesimal f\/lag structure with the regular, normal parabolic geometry inducing it.

In \cite{CS03}, \v{C}ap and Slov\'ak def\/ine a \emph{Weyl structure} for any parabolic geometry $(\Gbdle \rightarrow M,\omega)$ of type $(G,P)$, to be a $G_0$-equivariant section $\sigma: \Gbdle_0 \rightarrow \Gbdle$ of the $P_+$-principal bundle $\pi_+:\Gbdle \rightarrow \Gbdle_0$. We denote the set of Weyl structures by $\mathrm{Weyl}(\Gbdle,\omega)$. By Proposition~3.2 of~\cite{CS03}, global Weyl structures always exist for parabolic geometries in the real (smooth) category, and they exist locally in the holomorphic category. Considering the pull-back of the Cartan connection, $\sigma^*\omega$, the $\vert k \vert$-grading of $\g$ gives a $G_0$-invariant decomposition into components, $\sigma^*\omega = \sigma^*\omega_{-k} + \cdots + \sigma^*\omega_k,$ and by the observation that $\sigma$ commutes with fundamental vector f\/ields (i.e.\ $T_u\sigma(\widetilde{X}(u)) = \widetilde{X}(\sigma(u))$ for $X \in \g_0$ and $\widetilde{X}$ denoting the fundamental vector f\/ields of $X$ on $\Gbdle_0$ and $\Gbdle$) and the def\/ining properties of the Cartan connection, it follows that $\sigma^*\omega_i$ is horizontal for all $i \neq 0$, and that $\sigma^*\omega_0$ def\/ines a principal $G_0$ connection on $\Gbdle_0 \rightarrow M$ (cf.~3.3 of~\cite{CS03}). In particular, we see that the pair $(\Gbdle_0 \rightarrow M,\sigma^*\omega_{\leq})$ def\/ines a Cartan geometry of type $(P^*,G_0)$, where $P^* \isom \mathrm{exp}(\g_-) \rtimes G_0$ is the subgroup of $G$ containing $G_0$ with Lie algebra $\p^* = \g_- \dsum \g_0$, and the Cartan connection is given by
\[
\sigma^*\omega_{\leq} = \sigma^*\omega_{-k} + \cdots + \sigma^*\omega_0 \in \Omega^1(\Gbdle_0,\p^*).
\]

One reason Weyl structures are very useful for studying a parabolic geometry, is that they are in fact determined by very simple induced geometric objects, namely by the $\R_+$-principal connections they induce on certain ray bundles associated to $\Gbdle_0$. Fix an element $E_{\lambda}$ in the center of the reductive Lie algebra $\g_0$ such that $\mathrm{ad}(E_{\lambda})$ acts by scalar multiplication on each grading component $\g_i$ of $\g$ (for example the grading element $E$, which always exists and satisf\/ies $\mathrm{ad}(E)_{\vert \g_i} = i \cdot$). Then there is a unique representation $\lambda: G_0 \rightarrow \R_+$ satisfying $\lambda'(A) = B(E_{\lambda},A)$ for all $A \in \g_0$, $B$ the Killing form, and hence an associated $\R_+$-principal bundle $\mathcal{L}^{\lambda} \rightarrow M$. We have $\mathcal{L}^{\lambda} \isom \Gbdle_0/\mathrm{Ker}(\lambda)$ so let us denote the projection $\pi_{\lambda}: \Gbdle_0 \rightarrow \mathcal{L}^{\lambda}$. For any Weyl structure $\sigma$, the $1$-form $\lambda' \circ \sigma^*\omega_0 \in \Omega^1(\Gbdle_0)$ induces a $\R_+$-principal connection $\sigma^{\lambda}$ on $\mathcal{L}^{\lambda}$. After introducing these objects and studying their properties in Section~3 of~\cite{CS03}, \v{C}ap and Slov\'ak prove the fundamental result that the correspondence $\sigma \mapsto \sigma^{\lambda}$ def\/ines a bijective correspondence between the set of Weyl structures and the set of principal connections on $\mathcal{L}^{\lambda}$ (cf.\ Theorem~3.12 of \cite{CS03}).

{\sloppy In particular, this fact makes it possible to def\/ine certain distinguished classes of Weyl structures: A Weyl structure $\sigma$ is \emph{closed} if the induced $\R_+$-principal connection $\sigma^{\lambda}$ has vanishing curvature; it is \emph{exact} if $\sigma^{\lambda}$ is a trivial connection induced by a global trivialisation of the scale bundle $\mathcal{L}^{\lambda} \rightarrow M$. \v{C}ap and Slov\'ak prove that closed and exact Weyl structures always exist (in the smooth category), and the spaces of closed and exact Weyl structures are af\/f\/ine spaces over the closed, respectively over the exact, $1$-forms on $M$. Assuming a scale bundle $\mathcal{L}^{\lambda}$ to be f\/ixed, we denote the set of exact Weyl structures by $\overline{\mathrm{Weyl}}(\Gbdle,\omega)$, which is naturally identif\/ied with the set of global sections of $\mathcal{L}^{\lambda}$, and note that this set is non-empty (cf.\ Proposition~3.7 of~\cite{CS03}). Equivalently, an exact Weyl structure $\sigma$ is characterized by the existence of a holonomy reduction of the $G_0$-principal connection $\sigma^*\omega_0$ to the subgroup $\mathrm{Ker}(\lambda) \subset G_0$ (cf.\ Sections 3.13, 3.14 of~\cite{CS03}). We will denote this reduction by $r: \overline{\Gbdle}_0 \hookrightarrow \Gbdle_0$, and the corresponding reduction of~$\Gbdle$ to the structure group $\mathrm{Ker}(\lambda)$ by
\[
\overline{\sigma} := \sigma \circ r: \ \overline{\Gbdle}_0 \rightarrow \Gbdle.
\]
Thus an exact Weyl structure determines a Cartan geometry $(\overline{\Gbdle}_0 \rightarrow M,\overline{\sigma}^*\omega_{\leq})$ of type $(\overline{P^*},\mathrm{Ker}(\lambda))$ for $\overline{P^*} \isom \mathrm{exp}(\g_-) \rtimes \mathrm{Ker}(\lambda)$ the subgroup of $G$ containing $\mathrm{Ker}(\lambda)$ with Lie algebra $\overline{\p^*} := \g_- \dsum \mathrm{Ker}(\lambda')$. In the conformal case, exact Weyl structures correspond to metrics in the conformal equivalence class, while a general Weyl structure is given by a Weyl connection, i.e.\ a~torsion free connection which preserves the conformal equivalence class.

}

\subsection{Def\/inition and basic properties of essential structures}\label{section2.2}

Now we are ready to def\/ine essential parabolic structures and essential inf\/initesimal automorphisms. For now, let us take the following def\/initions for automorphisms, respectively inf\/initesimal automorphisms, of a Cartan geometry. For a Cartan geometry $(\Gbdle \rightarrow M,\omega)$ of arbitrary type $(G,P)$, an \emph{automorphism} $\Phi \in \mathrm{Aut}(\Gbdle,\omega)$ is a $P$-principal bundle morphism of~$\Gbdle$ such that $\Phi^*\omega = \omega$. An \emph{infinitesimal automorphism} $\mathbf{X} \in \mathrm{inf}(\Gbdle,\omega)$ is given by $\mathbf{X} \in \VF(\Gbdle)$, such that $(R_p)_*\mathbf{X} = \mathbf{X}$ and the Lie derivative satisf\/ies $\mathcal{L}_{\mathbf{X}} \omega = 0$.

Note that when $(\Gbdle,\omega)$ is a parabolic geometry, we get naturally induced $G_0$-bundle morphisms $\Phi_0:\Gbdle_0 \rightarrow \Gbdle_0$, $\R_+$-bundle morphisms $\Phi_{\lambda}: \mathcal{L}^{\lambda} \rightarrow \mathcal{L}^{\lambda}$, and dif\/feomorphisms $\varphi: M \rightarrow M$ for the $\Phi \in \mathrm{Aut}(\Gbdle,\omega)$. These are induced thanks to the $P$-equivariance of $\Phi$, by using commutativity with the appropriate projections, i.e.\ the def\/ining identities are $\Phi_0 \circ \pi_+ = \pi_+ \circ \Phi$, $\Phi_{\lambda} \circ \pi_{\lambda} = \pi_{\lambda} \circ \Phi_0$  and $\varphi \circ \pi = \pi \circ \Phi$. Hence, we get a natural action of the automorphism group $\mathrm{Aut}(\Gbdle,\omega)$ on the set of (exact) Weyl structures, by def\/ining, for $\Phi \in \mathrm{Aut}(\Gbdle,\omega)$, $\sigma \in \mathrm{Weyl}(\Gbdle,\omega)$,
\[
\Phi^*\sigma := \Phi^{-1} \circ \sigma \circ \Phi_0: \ \Gbdle_0 \rightarrow \Gbdle,
\]
and for $\sigma \in \overline{\mathrm{Weyl}}(\Gbdle,\omega)$ with corresponding global scale $s_{\sigma} \in \Gamma(\mathcal{L}^{\lambda})$,
\[
\Phi^*s_{\sigma} := \Phi_{\lambda}^{-1} \circ s_{\sigma} \circ \varphi: \ M \rightarrow \mathcal{L}^{\lambda}.
\]
Using the def\/ining relation $\pi_+ \circ \Phi^{-1} = \Phi_0^{-1} \circ \pi_+$ for $\Phi_0^{-1}$ and the corresponding relation bet\-ween~$\Phi_{\lambda}^{-1}$ and $\varphi^{-1}$, we can verify that $\Phi^*\sigma \in \Gamma(\Gbdle \rightarrow \Gbdle_0)$ and $\Phi^*s_{\sigma} \in \Gamma(\mathcal{L}^{\lambda} \rightarrow M)$; furthermore, $\Phi^*\sigma$ is $G_0$-equivariant by the equivariance properties of $\sigma$, $\Phi_0$ and $\Phi^{-1}$; i.e.\ $\Phi^*\sigma \in \mathrm{Weyl}(\Gbdle,\omega)$ and $\Phi^*s_{\sigma} \in \overline{\mathrm{Weyl}}(\Gbdle,\omega)$.

Similarly, for an inf\/initesimal automorphism $\mathbf{X} \in \VF(\Gbdle)$, we get naturally induced vector f\/ields $\mathbf{X}_0 \in \VF(\Gbdle_0)$, $\mathbf{X}_{\lambda} \in \VF(\mathcal{L}^{\lambda})$ and $X \in \VF(M)$ (with the f\/irst two being invariant with respect to the appropriate right-actions). For an arbitrary point of $M$, we may choose $\varepsilon > 0$ suf\/f\/iciently small so that $\Phi_{\mathbf{X},t}$, $\Phi_{\mathbf{X},t}^{-1}$, $\Phi_{\mathbf{X}_0,t}$, etc. all exist for $-\varepsilon \leq t \leq \varepsilon$, and so on this interval we have a~well-def\/ined family
\[
\Phi_{\mathbf{X},t}^*\sigma := \Phi_{\mathbf{X},t}^{-1} \circ \sigma \circ \Phi_{\mathbf{X}_0,t} \in \mathrm{Weyl}(\Gbdle,\omega),
\]
for any $\sigma \in \mathrm{Weyl}(\Gbdle,\omega)$, which is dif\/ferentiable in $t$ at $t=0$, so we can def\/ine the Lie derivative $\mathcal{L}_{\mathbf{X}} \sigma := (d/dt)\vert_{t=0}\Phi_{\mathbf{X},t}^*\sigma$; this is an element of the vector space on which the space of Weyl sections (an af\/f\/ine space) is modeled, i.e.\ it can be thought of as a $1$-form on $M$. For $s_{\sigma} \in \overline{\mathrm{Weyl}}(\Gbdle,\omega)$, the Lie derivative $\mathcal{L}_{\mathbf{X}}s_{\sigma}$ is def\/ined analogously, and it can be identif\/ied with an exact $1$-form on $M$.

\begin{Definition} \label{parabolic essential definition} For $(\Gbdle \rightarrow M,\omega)$ a parabolic geometry of type $(G,P)$, $\Phi \in \mathrm{Aut}(\Gbdle,\omega)$ an automorphism of the geometry, and $\sigma \in \mathrm{Weyl}(\Gbdle,\omega)$ a Weyl structure, $\Phi$ is an \emph{automorphism of~$\sigma$}, written $\Phi \in \mathrm{Aut}(\sigma)$, if and only if $\Phi^*\sigma = \sigma$ (equivalently, $\Phi \circ \sigma = \sigma \circ \Phi_0$). If $\sigma \in \overline{\mathrm{Weyl}}(\Gbdle,\omega)$ is an exact Weyl structure corresponding to $s_{\sigma} \in \Gamma(\mathcal{L}^{\lambda})$, then $\Phi$ is an \emph{exact automorphism of $\sigma$}, written $\Phi \in \mathrm{Aut}(\overline{\sigma})$, if and only if $\Phi^*s_{\sigma} = s_{\sigma}$ (equivalently, $\Phi_{\lambda} \circ s_{\sigma} = s_{\sigma} \circ \varphi$). An automorphism $\Phi \in \mathrm{Aut}(\Gbdle,\omega)$ is \emph{essential} if it is not an exact automorphism of any exact Weyl structure $\sigma \in \overline{\mathrm{Weyl}}(\Gbdle,\omega)$. We call $(\Gbdle,\omega)$ an \emph{essential parabolic structure} if $\mathrm{Aut}(\overline{\sigma}) \subsetneqq \mathrm{Aut}(\Gbdle,\omega)$ for every exact Weyl structure $\sigma$. We call a regular inf\/initesimal f\/lag structure $\mathcal{M} = (M,\{T^iM\},\Gbdle_0)$ of type $(\g,P)$ an \emph{essential structure} if the regular, normal parabolic geometry inducing it is essential.


An inf\/initesimal automorphism $\mathbf{X} \in \mathrm{inf}(\Gbdle,\omega)$ is an \emph{infinitesimal automorphism of $\sigma$}, written $\mathbf{X} \in \mathrm{inf}(\sigma)$, if and only if $\mathcal{L}_{\mathbf{X}} \sigma = 0$. For an exact Weyl structure $\sigma \in \overline{\mathrm{Weyl}}(\Gbdle,\omega)$, $\mathbf{X}$ is \emph{an exact infinitesimal automorphism of $\sigma$}, written $\mathbf{X} \in \mathrm{inf}(\overline{\sigma})$, if and only if $\mathcal{L}_{\mathbf{X}} s_{\sigma} = 0$. An inf\/initesimal automorphism is \emph{essential} if it is not an exact inf\/initesimal automorphism for any exact Weyl structure. For $\mathcal{M}$ a regular inf\/initesimal f\/lag structure as above, we say a vector f\/ield $X \in \VF(M)$ is an \emph{essential infinitesimal automorphism} of $\mathcal{M}$ if it lifts to an essential inf\/initesimal automorphism of the canonical parabolic geometry inducing $\mathcal{M}$.
\end{Definition}

\begin{Remark} Charles Frances has pointed out to us the def\/inition of essential parabolic structure given in Section~\ref{section2.2} of \cite{FrEss}, which did not make explicit use of Weyl structures. That def\/inition turns out to be almost equivalent to the one above, cf.\ Lemma \ref{inessential automorphisms TFAE}, except that in some cases the semi-simple part of the reductive group $G_0$ can be properly contained in $\mathrm{Ker}(\lambda)$. For example, in the case of strictly pseudoconvex CR structures, $G_0 \isom \R_+ \times U(n)$ and $\mathrm{Ker}(\lambda) \isom U(n)$ for an appropriate choice of scale $\lambda$. An exact Weyl structure is seen to be equivalent to a choice of pseudo-hermitian form for the CR structure, and so the above def\/inition of exact structure amounts to requiring that the group of CR transformations preserving any pseudo-hermitian form is always properly contained in the group of CR transformations.
\end{Remark}

\begin{Remark} Def\/inition \ref{parabolic essential definition} recovers the classical def\/inition of essentiality when the regular inf\/initesimal f\/lag structure is given by a conformal semi-Riemannian structure $(M,c)$ of signature $(p,q)$. In that case, $G_0$ is just the conformal group $\R_+ \times O(p,q)$, $\Gbdle_0$ is the bundle of frames which are semi-orthonormal with respect to some metric $g \in c$, and the choice of scale representation,
\begin{gather*}
 \lambda: \ \R_+ \times O(p,q) \rightarrow \R_+, \qquad
 \lambda: \ (s,A) \mapsto s^{-1},
\end{gather*}
identif\/ies $\mathcal{L}^{\lambda} \isom \Gbdle_0/\mathrm{Ker}(\lambda)$ with the ray bundle $\mathcal{Q} \rightarrow M$ of metrics in the conformal class, with the standard $\R_+$-action given by $g_x.s := s^2g_x$ for any $g \in c$ and $x \in M$ corresponding to $g_x \in \mathcal{Q}$. Exact Weyl structures thus correspond to choices of a metric in the conformal class, and a~conformal dif\/feomorphism $\varphi$ which uniquely corresponds to an automorphism $\Phi \in \mathrm{Aut}(\Gbdle,\omega)$ of the canonical conformal Cartan geometry is essential in the sense of Def\/inition~\ref{parabolic essential definition} if $\Phi \notin \mathrm{Aut}(\overline{\sigma})$ for all exact Weyl structures $\sigma$, i.e.\ if $\varphi$ fails to preserve all metrics in the conformal class.
\end{Remark}

\begin{Remark} \label{underlying automorphisms} Let $\mathcal{M} = (M,\{T^iM\},\Gbdle_0)$ denote a regular inf\/initesimal f\/lag structure of some parabolic type $(\g,P)$, if necessary including the extra geometric data required so that the regular, normal parabolic geometry of type $(G,P)$ inducing it is unique up to isomorphism. If we need to distinguish this parabolic geometry from others of the same type, we will use the notation $(\Gbdle,\omega^{nc})$ to signify the canonical (\textbf{n}ormal \textbf{C}artan) geometry. In this setting, we can def\/ine an automorphism of the structure in terms of $\mathcal{M}$: An automorphism of the regular inf\/initesimal f\/lag structure, $\varphi \in \mathrm{Aut}(\mathcal{M})$, is a dif\/feomorphism $\varphi \in \mathrm{Dif\/f}(M)$ which satisf\/ies: $(i)$~$\varphi_*(T_x^iM) \subseteq T_{\varphi(x)}^iM$ for all $x \in M$ and all $-k \leq i \leq -1$; and $(ii)$~the induced bundle map $\mathrm{gr}(\varphi)$ (which as a~consequence of $(i)$~is a lift of $\varphi$ def\/ined on the bundle $\mathcal{F}(\mathrm{gr}(TM))$ of frames of the associated graded tangent bundle) preserves $\Gbdle_0$ as a subbundle of $\mathcal{F}(\mathrm{gr}(TM))$ (and hence $\mathrm{gr}(\varphi)$ restricts to a $G_0$-bundle morphism $\Phi_0$ of $\Gbdle_0$). We can identify $\mathrm{Aut}(\mathcal{M})$ with $\mathrm{Aut}(\Gbdle,\omega^{nc})$ since by uniqueness of $(\Gbdle,\omega^{nc})$ up to isomorphism, $\varphi$ (and $\Phi_0$) lift to a unique $P$-bundle morphism $\Phi$ of $\Gbdle$ preserving $\omega^{nc}$ under pullback. Thus, we can think of an automorphism $\varphi \in \mathrm{Aut}(\mathcal{M})$ as including as well the automorphism $\Phi \in \mathrm{Aut}(\Gbdle,\omega^{nc})$ and the induced $G_0$-bundle morphism $\Phi_0 = \mathrm{gr}(\varphi)_{\vert \Gbdle_0}$ of $\Gbdle_0$.
\end{Remark}

\begin{Lemma} \label{inessential automorphisms TFAE} Let $\Phi \in \mathrm{Aut}(\Gbdle,\omega)$ be an automorphism of a parabolic geometry, let $\sigma$ be a Weyl structure, and let a scale bundle $\mathcal{L}^{\lambda} \rightarrow M$ be fixed. The following are equivalent:
\begin{enumerate}\itemsep=0pt
\item[$(i)$] $\Phi \in \mathrm{Aut}(\sigma)$;
\item[$(ii)$] For the induced bundle morphism $\Phi_0: \Gbdle_0 \rightarrow \Gbdle_0$, we have $\Phi_0 \in \mathrm{Aut}(\Gbdle_0,\sigma^*\omega_{\leq})$;
\item[$(iii)$] $\Phi_{\lambda}$ preserves the scale bundle connection $\sigma^{\lambda} \in \Omega^1(\mathcal{L}^{\lambda})$: $\Phi_{\lambda}^*\sigma^{\lambda} = \sigma^{\lambda}$.
\end{enumerate}

 If $\sigma$ is exact, then $\Phi \in \mathrm{Aut}(\overline{\sigma})$ if and only if the induced bundle morphism $\Phi_0$ preserves the sub-bundle $\overline{\Gbdle}_0 \subset \Gbdle_0$ and the restriction satisfies $(\Phi_0)_{\vert \overline{\Gbdle}_0} \in \mathrm{Aut}(\overline{\Gbdle}_0,\overline{\sigma}^*\omega_{\leq})$.

For $\mathbf{X} \in \mathrm{inf}(\Gbdle,\omega)$, $\mathbf{X} \in \mathrm{inf}(\sigma)$ if and only if $\mathbf{X}_0 \in \mathrm{inf}(\Gbdle_0,\sigma^*\omega_{\leq})$ and, for $\sigma$ exact, $\mathbf{X} \in \mathrm{inf}(\overline{\sigma})$ if and only if the restriction of $\mathbf{X}_0$ to $\overline{\Gbdle}_0$ is tangent to $\overline{\Gbdle}_0$ and induces an element of $\mathrm{inf}(\overline{\Gbdle}_0,\overline{\sigma}^*\omega_{\leq})$.
\end{Lemma}

\begin{proof} $(i)$ $\Rightarrow$ $(ii)$, since by $\Phi \circ \sigma = \sigma \circ \Phi_0$ and $\Phi \in \mathrm{Aut}(\Gbdle,\omega)$ we have $\Phi_0^*(\sigma^*\omega) = \sigma^*\omega$, and in particular $\Phi_0^*(\sigma^*\omega_{\leq}) = \sigma^*\omega_{\leq}$. And $(ii)$ $\Rightarrow$ $(iii)$, since $\Phi_0^*(\sigma^*\omega_{\leq}) = \sigma^*\omega_{\leq}$ implies that $\Phi_0^*(\sigma^*\omega_0) = \sigma^*\omega_0$. In particular, $\Phi_0^*(\lambda' \circ \sigma^*\omega_0) = \lambda' \circ \Phi_0^*(\sigma^*\omega_0) = \lambda' \circ \sigma^*\omega_0$. Hence, the $\R_+$-bundle morphism $\Phi_{\lambda}$ and the $\R_+$-principal connection $\sigma^{\lambda}$, induced on $\mathcal{L}^{\lambda}$ by $\Phi_0$ and $\lambda' \circ \sigma^*\omega_0$, respectively, satisfy: $\Phi_{\lambda}^* \sigma^{\lambda} = \sigma^{\lambda}$.

  We now show $(iii)$ $\Rightarrow$ $(i)$: Consider the $\R_+$-principal connection $(\Phi^*\sigma)^{\lambda} \in \Omega^1(\mathcal{L}^{\lambda})$. It is induced by:
\begin{gather*}
\lambda' \circ (\Phi^*\sigma)^*\omega_0  = \lambda' \circ \big(\big(\Phi^{-1}\circ \sigma \circ \Phi_0\big)^*\omega_0\big)
  = \lambda' \circ \big(\Phi_0^*\sigma^*\big(\Phi^{-1}\big)^*\omega_0\big) \\
\phantom{\lambda' \circ (\Phi^*\sigma)^*\omega_0}{} = \lambda' \circ (\Phi_0^*\sigma^*\omega_0) = \Phi_0^*(\lambda' \circ \sigma^*\omega_0).
\end{gather*}
So $(\Phi^*\sigma)^{\lambda} = \Phi_{\lambda}^*\sigma^{\lambda}$, which equals $\sigma^{\lambda}$ by assumption $(iii)$. Thus, by Theorem~3.12 of \cite{CS03} (cf.\ discussion in Section~\ref{section2.1}), the Weyl structures $\Phi^*\sigma$ and $\sigma$ are equal, showing that $(i)$ holds.

 To see the f\/inal statement of the lemma, let $\sigma$ be an exact Weyl structure and let us denote by $s_{\sigma} \in \Gamma(\mathcal{L}^{\lambda})$ the global scale which induces the trivial connection $\sigma^{\lambda} \in \Omega^1(\mathcal{L}^{\lambda})$. That is, for any point $p = s_{\sigma}(x).r \in \mathcal{L}^{\lambda}$, for $x \in M$, $r \in \R_+$, we have the decomposition
\[
T_p \mathcal{L}^{\lambda} = (R_r)_*((s_{\sigma})_*(T_xM)) \dsum \R \zeta_1(p),
\]
for $\zeta_1$ the fundamental vector f\/ield on $\mathcal{L}^{\lambda}$ of the vector $1 \in \R$; the value of $\sigma^{\lambda}$ on a tangent vector $v \in T_p\mathcal{L}^{\lambda}$ is given by the coef\/f\/icient of $\zeta_1(p)$ determined by this decomposition. Then the holonomy reduction of $(\mathcal{L}^{\lambda},\sigma^{\lambda})$ to the trivial structure group is given by $s_{\sigma}(M) \subset \mathcal{L}^{\lambda}$, and the reduction of $(\Gbdle_0,\sigma^*\omega_0)$ to $\mathrm{Ker}(\lambda)$ is given by
\[
\overline{\Gbdle}_0 = \pi_{\lambda}^{-1}(s_{\sigma}(M)) \subset \Gbdle_0.
\]
Hence for $x \in M$ and $u \in (\overline{\Gbdle}_0)_x$, we have $\pi_{\lambda}(u) = s_{\sigma}(x)$ and $\Phi_0(u) \in \overline{\Gbdle}_0$ if and only if $\pi_{\lambda}(\Phi_0(u)) = s_{\sigma}(\pi_0(\Phi_0(u)))$. But $\pi_{\lambda}(\Phi_0(u)) = \Phi_{\lambda}(\pi_{\lambda}(u)) = \Phi_{\lambda}(s_{\sigma}(x))$ and $\pi_0(\Phi_0(u)) = \varphi(\pi_0(u)) = \varphi(x)$. Thus, $\Phi_0(u) \in \overline{\Gbdle}_0$ if and only if $\Phi_{\lambda}(s_{\sigma}(x)) = s_{\sigma}(\varphi(x))$. But if $\Phi \in \mathrm{Aut}(\overline{\sigma})$, then clearly $\Phi_{\lambda}$ preserves the induced (trivial) connection $\sigma^{\lambda} \in \Omega^1(\mathcal{L}^{\lambda})$, so by $(iii)$ $\Leftrightarrow$ $(ii)$ shown above we have $\Phi_0^*(\sigma^*\omega_{\leq}) = \sigma^*\omega_{\leq}$ and since $\Phi_0$ preserves $\overline{\Gbdle}_0$ the corresponding identity follows for the restriction and $\overline{\sigma}^*\omega_{\leq}$.

 The statements for $\mathbf{X} \in \mathrm{inf}(\Gbdle,\omega)$ are proven in the same manner.
 \end{proof}

\begin{Remark} In particular, it follows from the proof of Lemma \ref{inessential automorphisms TFAE} that any exact automor\-phism~$\Phi$ of an exact Weyl structure $\sigma$ is also an automorphism of $\sigma$, i.e.\ $\Phi \in \mathrm{Aut}(\overline{\sigma}) \Rightarrow \Phi \in \mathrm{Aut}(\sigma)$ as one would hope.

 On the other hand, the converse does not hold: If $\Phi \in \mathrm{Aut}(\sigma)$ for an exact Weyl structure $\sigma$, the requirement that $\Phi_0(\overline{\Gbdle}_0) \subset \overline{\Gbdle}_0$ is necessary to guarantee that an automorphism $\Phi$ of an exact Weyl structure $\sigma$ is in fact an \emph{exact} automorphism. An instructive example is the conformal structure induced by the Euclidean metric on $\R^n$: The dif\/feomorphism given by dilation by a~positive constant $r$ is always an automorphism of the exact Weyl structure corresponding to the Euclidean metric. However, for $r \neq 1$, this dif\/feomorphism is not an isometry, hence not an \emph{exact} automorphism. We are grateful to Felipe Leitner for bringing this to our attention, which led us to modify an earlier version of Def\/inition~\ref{parabolic essential definition}.
\end{Remark}
